
\section{Subconjuntos fechados}
Ao longo desta seção, $(M,d)$ denotará um espaço métrico arbitrário.

Intuitivamente, um conjunto fechado é um conjunto que contém todos os pontos do espaço ambiente que estão ``grudados'' nele.

\begin{definition}\index{conjunto fechado}
    Seja $A$ um subconjunto de $M$.
    Dizemos que $A$ é \emph{fechado} se, para todo ponto de acumulação $p$ de $A$, temos que $p \in A$.
\end{definition}

Frequentemente, uma forma mais prática de verificar se um conjunto é fechado é dada pelo seguinte critério.
\begin{proposition}\label{prop:fechado-complemento-aberto}
    Seja $A$ um subconjunto de $M$.
    Então, $A$ é fechado se, e somente se, seu complemento $M\setminus A$ é aberto.
\end{proposition}

\begin{proof}
    Primeiramente, suponha que $A$ é fechado.
    Veremos que $M\setminus A$ é aberto.

    Seja $p \in M\setminus A$ arbitrário.
    O ponto $p$ não é um ponto de acumulação de $A$, pois, caso contrário, como $A$ é fechado, teríamos que $p \in A$, o que é uma contradição.

    Assim, como $p$ não é ponto de acumulação de $A$, existe $r>0$ tal que $\bigl(B(p, r)\cap A\bigr)\setminus\{p\} = \emptyset$.
    Daí, $B(p, r)\cap A\subseteq \{p\}$.
    Como $p \notin A$, temos que $B(p, r)\cap A = \emptyset$.
    Dessa forma, $B(p, r)\subseteq M\setminus A$.

    Reciprocamente, suponha que $M\setminus A$ é aberto.
    Veremos que $A$ é fechado.
    Para tanto, devemos ver que para todo $p\in M$, se $p$ é um ponto de acumulação de $A$, então $p \in A$.
    Pela contrapositiva, basta ver que para todo $p\in M$, se $p \notin A$, então $p$ não é um ponto de acumulação de $A$.

    De fato, dado $p \in M\setminus A$, como $M\setminus A$ é aberto, existe $r>0$ tal que $B(p, r)\subseteq M\setminus A$, de modo que $B(p, r)\cap A = \emptyset$.
    Assim, $\bigl(B(p, r)\cap A\bigr)\setminus\{p\} = \emptyset$, o que mostra que $p$ não é um ponto de acumulação de $A$.
\end{proof}

\begin{example}
    O intervalo $[a, b]$ com $a<b$ é um conjunto fechado de $\mathbb R$, pois $\mathbb R\setminus [a, b] = ]-\infty, a[ \cup ]b, +\infty[$ é uma união de conjuntos abertos, logo é aberto.
\end{example}

\begin{corollary}
    Seja $A$ um subconjunto de $M$.
    Então, $A$ é aberto se, e somente se, seu complemento $M\setminus A$ é fechado.
\end{corollary}
\begin{proof}
    Seja $B=M\setminus A$.
    Pela Proposição~\ref{prop:fechado-complemento-aberto}, $B$ é fechado se, e somente se, $M\setminus B$ é aberto.
    Como $M\setminus B=M\setminus(M\setminus A)=A$, segue que
    \[
        M\setminus A \text{ é fechado}
        \quad\Longleftrightarrow\quad
        A \text{ é aberto}.
    \]
    Portanto, $A$ é aberto se, e somente se, $M\setminus A$ é fechado.
\end{proof}

\begin{proposition}
    Seja $F$ um subconjunto finito de $M$.
    Então $F$ não possui pontos de acumulação.
    Em particular, $F$ é fechado.
\end{proposition}
\begin{proof}
    Seja $p\in M$ arbitrário.
    Para cada $a\in F\setminus\{p\}$, seja $r_a=d(p,a)>0$.

    Se $F\setminus\{p\}=\emptyset$, tome $r=1$.
    Caso contrário, como $F\setminus\{p\}$ é finito e não vazio, tome
    \[
        r=\min\{r_a:a\in F\setminus\{p\}\}.
    \]

    Em ambos os casos, temos $r>0$ e $\bigl(B(p,r)\cap F\bigr)\setminus\{p\}=\emptyset$.
    Logo, $p$ não é ponto de acumulação de $F$.

    Como $p\in M$ foi arbitrário, $F$ não possui pontos de acumulação.
    Portanto, por vacuidade, $F$ é fechado.
\end{proof}

\begin{proposition}
    Sejam $A$ um subconjunto aberto de $M$ e $C$ um subconjunto fechado de $M$.
    Então $A\setminus C$ é um conjunto aberto de $M$.
\end{proposition}
\begin{proof}
    Como $C$ é fechado, segue da Proposição~\ref{prop:fechado-complemento-aberto} que $M\setminus C$ é aberto.
    Assim,
    \[
        A\setminus C = A\cap(M\setminus C)
    \]
    é uma interseção de dois conjuntos abertos.
    Logo, $A\setminus C$ é aberto.
\end{proof}